\documentclass[11pt,a4paper]{article}
\usepackage[margin=23mm]{geometry}
\usepackage{amsmath,amssymb,amsthm,hyperref}
\newtheorem{theorem}{Theorem}
\newtheorem{lemma}[theorem]{Lemma}
\newtheorem{proposition}[theorem]{Proposition}
\newcommand{\Om}{\Omega}
\newcommand{\om}{\omega}
\newcommand{\dd}{\,d}
\title{Repeated prime factors and strong three-dimensional partitions}
\author{Human Compression project: supplementary proof for F008}
\date{}
\begin{document}\maketitle
Theorem 1.3 of Lebowitz--Lockard and Tenenbaum~\cite{LLT} concerns
integral partitions of a cube into equal-volume boxes whose lengths
along one fixed axis are pairwise distinct. We prove that the set of
side lengths admitting such a partition has an asymptotic density
strictly less than one. This supplements the already reviewed complete
companions \emph{Divisor concentration and primitive bases} and
\emph{Primitive rectangular partitions and planar density}. We acquire
both companions in full; no part is treated as a free earlier argument.

The source's primitive-basis and prime-divisor strategy is retained and
attributed. The new argument replaces the divisor decomposition and
Fourier estimates in its Section 6. In particular, we impose a stronger
small-prime cutoff and assign each divisor its \emph{largest repeated
large prime}. That assignment makes the two-prime constraints
independent. No finite calculation is a premise of this uniform proof.

Write $\log_2x=\log\log x$. All constants below are independent of
$x$, which is sufficiently large. We use Chebyshev's bound
$\pi(u)\ll u/\log u$ and Euler summation for the zeta function,
already retained in the companions. The parameters in the final
application are fixed as
\[
 \varepsilon=\frac1{100},\qquad c=\frac{49}{50},\qquad
 \kappa=\frac{21}{20}.
\]

\section{A cutoff that forces a repeated large prime}
Put $L=\log x$, $\xi=\log L$, and
\[
 z=\exp(L/\xi^2),\qquad Y=\exp(L/(10\xi)),\qquad
 n_z=\prod_{\substack{\ell^a\parallel n\\\ell\le z}}\ell^a.
\]
Here and below letters $p,q,\ell$ used as prime factors denote primes.

\begin{lemma}[Small-prime part]\label{smallpart}
For every fixed $A>0$,
\[
 \#\{n\le x:n_z>Y\}\ll_A xL^{-A}.
\]
\end{lemma}
\begin{proof}
Set $\lambda=\log\xi/\log z$; eventually $0<\lambda<1/4$.
The multiplicative function $f(n)=n_z^\lambda$ has nonnegative
divisor-sum coefficients: at $\ell\le z$ these are
$g(\ell^a)=\ell^{a\lambda}-\ell^{(a-1)\lambda}$ for $a\ge1$,
and at $\ell>z$ they are zero. Consequently
\[
 \sum_{n\le x} f(n)
 \le x\sum_{d\ge1}\frac{g(d)}d
 =x\prod_{\ell\le z}\frac{1-1/\ell}{1-\ell^{\lambda-1}}.
\]
The logarithm of this product is
\[
 \sum_{\ell\le z}\frac{\ell^\lambda-1}{\ell}+O(1),
\]
uniformly for $\lambda<1/4$, since all terms of degree at least two
are bounded by a constant times $\sum_\ell\ell^{-3/2}$.
Chebyshev's bound and partial summation give
\[
 \sum_{\ell\le z}\frac{\ell^\lambda-1}{\ell}
 \ll 1+\frac{z^\lambda}{\log z}
 +\int_2^z\frac{(t^\lambda-1)+\lambda t^\lambda}{t\log t}\dd t
 \ll 1+z^\lambda.
\]
For the last estimate substitute $v=\lambda\log t$.
On $v\le1$, $(e^v-1)/v$ is bounded and
$\lambda\int_{\lambda\log2}^1v^{-1}\dd v=O(1)$;
on $v\ge1$ both integrals are $O(e^{\lambda\log z})$.
As $z^\lambda=\xi$, Markov's inequality now gives
\[
 \#\{n\le x:n_z>Y\}
 \le x\exp\{-\lambda\log Y+O(\xi)\}
 =x\exp\{-\tfrac1{10}\xi\log\xi+O(\xi)\}.
\]
This is $O_A(xL^{-A})$ for every fixed $A$.
\end{proof}

Call $b$ \emph{good} if
\begin{equation}\label{good}
 \sqrt x<b\le x,\quad
 \Om(b)\le(1+\varepsilon)\log_2 b,\quad
 b_z\le Y,\quad \ell^2\mid b\Longrightarrow\ell\le z.
\end{equation}
The last two conditions discard at most $O_A(xL^{-A})$ integers from
the first two: use Lemma~\ref{smallpart} and
$x\sum_{j>z}j^{-2}\ll x/z$.

\begin{lemma}[Canonical repeated prime]\label{assignment}
Let $b$ be good, $p=P^+(b)$, and $m=b/p$. Then $p\nmid m$,
\[
 p\ge x^{1/(2(1+\varepsilon)\xi)}>z,\qquad
 \tau(b^3)\le4^{\Om(b)}\le L^{(1+\varepsilon)\log4}<L^2.
\]
Every divisor $d\mid m^3$ in $[b/L^2,b)$ has a prime $q>z$
whose exponent in $d$ is two or three. Thus it has a well-defined
largest such prime, denoted $q(d)$. If $q(d_1)>q(d_2)$, the exponent
of $q(d_1)$ in $d_2$ is at most one.
\end{lemma}
\begin{proof}
Since $b\le p^{\Om(b)}$ and $b>\sqrt x$, the lower bound for $p$
follows. For large $x$ it exceeds $z$, so $p$ occurs to exponent one
by~\eqref{good}. The divisor bound follows from $3a+1\le4^a$.

If a divisor $d\mid m^3$ had exponent at most one at every prime
above $z$, then
\[
 d\le b_z^3\frac{m}{b_z}=b_z^2m
 \le\frac{Y^2b}{p}<\frac b{L^2}.
\]
The last inequality holds because
$\log(Y^2L^2)=L/(5\xi)+2\xi$
is smaller than $L/(2(1+\varepsilon)\xi)$.
This proves existence; the final assertion is the definition of the
largest repeated large prime.
\end{proof}

\section{A uniform weighted divisor-pair estimate}
For real $\theta$, put
$\tau(v,\theta)=\sum_{d\mid v}d^{i\theta}$.

\begin{lemma}[Fourier mean]\label{fourier}
For fixed $0<y\le1$, uniformly for $x\ge3$,
\[
 \int_0^1\sum_{h\le x}
       \frac{y^{\Om(h)}|\tau(h^3,\theta)|^2}{h}\dd\theta
 \ll_y L^{4y+(12y-1)^+}
       \begin{cases}1+\log L,&12y=1,\\1,&12y\ne1.\end{cases}
\]
In particular, at $y=1/12$ the bound is $O(L^{1/3}(1+\xi))$.
\end{lemma}
\begin{proof}
Take $\sigma=1+1/L$. Nonnegativity permits replacing the truncated
harmonic sum by at most $e$ times its full Dirichlet series
$F(\sigma,\theta)$. It is an Euler product, since
$y^{\Om(h)}|\tau(h^3,\theta)|^2$ is multiplicative.
Its terms at a prime power $\ell^a$ are at most
$y^a(3a+1)^2\ell^{-a\sigma}$; those with $a\ge2$ have uniformly
bounded total contribution to the logarithm.

The exact first-degree term is determined by
\begin{equation}\label{polynomial}
 |1+e^{iu}+e^{2iu}+e^{3iu}|^2
 =4+6\cos u+4\cos(2u)+2\cos(3u).
\end{equation}
Indeed each coefficient counts ordered exponent pairs of a specified
difference. Using the absolutely convergent logarithm of the Euler
product for $\zeta$, whose prime-power terms of degree at least two
are uniformly bounded, we obtain
\[
 F(\sigma,\theta)\ll_y
 \zeta(\sigma)^{4y}
 |\zeta(\sigma+i\theta)|^{6y}
 |\zeta(\sigma+2i\theta)|^{4y}
 |\zeta(\sigma+3i\theta)|^{2y}.
\]
This argument uses $\log(1+u)\le u$ on each positive local factor;
no assumption on zeros on the line $\Re s=1$ is required.

Euler summation gives, for $\Re s>1$,
\[
 \zeta(s)=\frac{s}{s-1}-s\int_1^\infty\{v\}v^{-s-1}\dd v.
\]
For $1<\Re s\le2$ and $|\Im s|\le3$ this implies
$|\zeta(s)|\ll1+|s-1|^{-1}$.
Therefore, for $0<\theta\le1$ and $j=1,2,3$,
$|\zeta(\sigma+ij\theta)|\ll\min(L,1/\theta)$, while
$\zeta(\sigma)\ll L$. Integrating
$L^{4y}\min(L,1/\theta)^{12y}$ over $0<\theta\le1$ proves
all three cases of the stated bound.
\end{proof}

\begin{lemma}[Shifted pair concentration]\label{pairs}
Let $a_1,a_2\in\{2,3\}$, $V\in\mathbb R$, and $W\ge3$.
Define
\[
 D(h)=\#\{(e_1,e_2):e_1,e_2\mid h^3,
       |a_1\log e_1-a_2\log e_2-V|\le W\}.
\]
Uniformly in $V$, for $W\ll\xi$ and $x$ large,
\[
 \sum_{\substack{h\le x\\\Om(h)\le(1+\varepsilon)\xi}}
       \frac{D(h)}h
 \ll \xi^2 L^{1/3+(1+\varepsilon)\log12}.
\]
\end{lemma}
\begin{proof}
On $|v|\le W$ the function
$\big(\sin(v/(2W))/(v/(2W))\big)^2$ is bounded below by a positive
absolute constant. Its Fourier representation is
\[
 \left(\frac{\sin(v/(2W))}{v/(2W)}\right)^2
 =W\int_{-1/W}^{1/W}(1-W|\theta|)e^{i\theta v}\dd\theta.
\]
Summing this majorant, taking absolute values, and using
$|AB|\le(|A|^2+|B|^2)/2$ show
\[
 D(h)\ll W\int_0^1|\tau(h^3,u)|^2\dd u.
\]
The substitutions $u=a_j\theta$ are valid because $a_j/W\le1$;
the phase $e^{-iV\theta}$ has modulus one.
For $y=1/12$, the restriction on $\Om(h)$ gives
$1\le y^{\Om(h)}L^{-(1+\varepsilon)\log y}$.
Apply Lemma~\ref{fourier}, and use $W\ll\xi$.
\end{proof}

\begin{lemma}[Prime windows]\label{primewindows}
For $z\ge3$ and a real interval $J$ of length $W\ge0$,
\[
 \sum_{\substack{q>z\\\log q\in J}}\frac1q
 \ll\frac{W+1}{\log z}.
\]
If primes $p>z$ are additionally restricted by $p\le T$, their number
is at most $O(T/\log z)$, with zero possibilities when $T\le z$.
\end{lemma}
\begin{proof}
Intersect the interval with $[\log z,\infty)$ and divide it into
at most $W+2$ intervals of length at most one. In such an interval
$[u,u+1]$, Chebyshev's bound gives
$\sum_{e^u<q\le e^{u+1}}q^{-1}\le e^{-u}\pi(e^{u+1})
\ll1/u\le1/\log z$.
The second assertion is Chebyshev's bound at $T>z$.
\end{proof}

\section{The two repeated-prime cases}
For good $b=mp$ put
\[
 T(b)=\#\{d\mid m^3:b/L^2\le d<b\},\qquad
 \alpha_1=(1+\varepsilon)\log12-\frac53,\quad
 \alpha_2=(1+\varepsilon)\log16-2.
\]

\begin{proposition}[Second moment near the largest prime]\label{second}
For $\sqrt x\le X\le x$,
\[
 \sum_{\substack{X<b\le\min(2X,x)\\b\ {
m good}}}T(b)^2
 \ll X\xi^8\big(L^{\alpha_1}+L^{\alpha_2}\big).
\]
\end{proposition}
\begin{proof}
Expand the square into ordered divisor pairs $(d_1,d_2)$ and use
the assignment of Lemma~\ref{assignment}. Every divisor in the sum
lies in $[X/L^2,2X)$.

\emph{Equal assigned primes.}
Suppose $q(d_1)=q(d_2)=q$. Write $b=hqp$ and
$d_j=e_jq^{\nu_j}$, where $e_j\mid h^3$ and
$\nu_j\in\{2,3\}$. This is legitimate because all primes above $z$
occur to exponent one in $b$, and $q<p$.
For fixed $h,q$, there are at most
$O(X/(hq\log z))$ possible primes $p$.
The restrictions on each divisor imply
\[
 \log X-2\xi\le\log e_j+\nu_j\log q<\log X+\log2.
\]
Their simultaneous validity implies
\[
 |\nu_2\log e_1-\nu_1\log e_2
       -(\nu_2-\nu_1)\log X|\le20\xi.
\]
For such a pair $(e_1,e_2)$, the allowed $\log q$ belong to an
interval of length $O(\xi)$, independently of $p$.
Lemma~\ref{primewindows} and then Lemma~\ref{pairs} give an upper
bound
\[
 \frac{X\xi}{(\log z)^2}
 \sum_{\substack{h\le x\\\Om(h)\le(1+\varepsilon)\xi}}
       \frac{D(h)}h
 \ll X\xi^7 L^{-2+1/3+(1+\varepsilon)\log12}
 =X\xi^7L^{\alpha_1}.
\]
There are only four choices for $(\nu_1,\nu_2)$. Coprimality and
ordering restrictions may safely be discarded in these upper bounds.
The shifted pair estimate is needed also when $\nu_1\ne\nu_2$;
ordinary unshifted closeness of $e_1,e_2$ is not assumed.

\emph{Distinct assigned primes.}
By interchanging the two divisors at a cost of at most two, suppose
$q_1=q(d_1)>q_2=q(d_2)$. Write $b=hq_1q_2p$ and
\[
 d_1=e_1q_1^{\nu_1}q_2^{\mu_1},\qquad
 d_2=e_2q_1^{\mu_2}q_2^{\nu_2},
\]
where $e_1,e_2\mid h^3$, $2\le\nu_1,\nu_2\le3$,
$0\le\mu_1\le3$, and $0\le\mu_2\le1$.
The last inequality is exactly the largest-prime rule. In particular
\[
 \det\begin{pmatrix}\nu_1&\mu_1\\\mu_2&\nu_2\end{pmatrix}
 =\nu_1\nu_2-\mu_1\mu_2\ge4-3=1.
\]
For fixed $h,e_1,e_2$ and these exponents, the two logarithmic divisor
restrictions define a parallelogram for
$(u,v)=(\log q_1,\log q_2)$. Each defining interval has length
$2\xi+\log2$. The inverse matrix has uniformly bounded coefficients,
so each coordinate projection of the parallelogram has length
$O(\xi)$, with endpoints independent of $p$.

For fixed $q_1,q_2$, there are
$O(X/(hq_1q_2\log z))$ possible $p$. Bounding the parallelogram
by its coordinate projections and using Lemma~\ref{primewindows}
twice gives the total upper bound
\[
 \frac{X\xi^2}{(\log z)^3}
 \sum_{\substack{h\le x\\\Om(h)\le(1+\varepsilon)\xi}}
       \frac{\tau(h^3)^2}{h}.
\]
Since $\tau(h^3)^2\le16^{\Om(h)}
\le L^{(1+\varepsilon)\log16}$ and
$\sum_{h\le x}h^{-1}\ll L$, this is
$O(X\xi^8 L^{(1+\varepsilon)\log16-2})$.
There are only finitely many exponent choices. The two cases exhaust
all ordered divisor pairs and prove the proposition.
\end{proof}

\begin{theorem}[Exceptional divisors]\label{exceptional}
For the fixed constants above,
\[
 \#\{b\ {
m good}:T(b)>L^c\}\ll xL^{-\kappa}.
\]
\end{theorem}
\begin{proof}
Apply Proposition~\ref{second} to dyadic intervals covering
$(\sqrt x,x]$. Their left endpoints have sum $O(x)$; thus there is
no additional factor $\log x$ in the resulting bound. Markov's
inequality gives
\[
 \#\{b\ {
m good}:T(b)>L^c\}
 \ll x\xi^8\big(L^{\alpha_1-2c}+L^{\alpha_2-2c}\big).
\]
For an entirely rational check of the strict margin, the exponential
series through degree eight gives $e^{5/2}>12$ and $e^{14/5}>16$.
Thus $\log12<5/2$, $\log16<14/5$, and
\[
 \alpha_1<\frac{103}{120},\qquad
 \alpha_2<\frac{207}{250}<\frac{103}{120},\qquad
 2c-\frac{103}{120}-\frac{21}{20}=\frac{31}{600}>0.
\]
For reference, the sharper values are
\[
 \alpha_1=0.8430890496\ldots,\qquad
 \alpha_2=0.8003146094\ldots,
\]
and $2c-\max(\alpha_1,\alpha_2)=1.1169109503\ldots$.
The fixed factor $\xi^8$ is absorbed by the rational strict margin
above; the decimals are explanatory.
\end{proof}

\section{The primitive basis and density}
A \emph{strong three-dimensional partition} has at least two integral
axis-parallel boxes of equal volume and has pairwise distinct lengths
along one specified axis. Let $B$ be the set of side lengths admitting
such a partition with no smaller divisor admitting one.

\begin{theorem}\label{density}
The set of strong three-dimensional Mondrian side lengths has an
asymptotic density strictly less than one. In particular, the conclusion
of source Theorem 1.3 holds.
\end{theorem}
\begin{proof}
We explicitly use three complete earlier arguments. The companion
common paper proves that the multiples of integers
$b$ with $\Om(b)>(1+\varepsilon)\log_2b$ have vanishing upper-density
tails. It also proves the finite Behrend inequality and the consequence
that a set of multiples has density less than one when its primitive
basis excludes one and has vanishing upper-density tails. The companion
application proves its geometric trace lemma and the ``few sums, few
prime factors'' counting lemma for bounded nonnegative coefficients.
These statements concern arbitrary dimension and are used with $k=3$.

Split $B=B_1\cup B_2$ according as the displayed prime-factor bound
fails or holds. The common tail lemma handles $B_1$. We show
$\#(B_2\cap[1,x])\ll xL^{-\kappa}$, so the reciprocal sum over
$B_2$ converges. Integers at most $\sqrt x$ are harmless. Among the
remaining integers of $B_2$, the nongood ones are $O_A(xL^{-A})$
by Lemma~\ref{smallpart} and the large prime-square bound.
Theorem~\ref{exceptional} discards $O(xL^{-\kappa})$ further
integers and leaves $T(b)\le L^c$.

For each remaining $b=mp$, the prime $p$ is a proper divisor: the
companion geometric lemma proves that a prime side length has no strong
partition. Primitivity therefore supplies a coordinate trace of length
$b$ with at least one segment not divisible by $p$.
Every segment length divides $b^3$, is at least
$b/\tau(b^3)\ge b/L^2$, and is less than $b$.
Let $D$ be the sum of the segments not divisible by $p$.
Then $1\le D\le b\le x$, and reducing the trace modulo $p$ gives
$p\mid D$. All its distinct summands divide $m^3$ and are among
the at most $T(b)\le L^c$ possibilities. Each multiplicity is at most
the total number of boxes, which is at most $\tau(b^3)\le L^2$.
The chosen trace need not be along the axis with distinct lengths;
this is precisely why bounded multiplicities are retained.

Moreover $p\ge x^{1/(2(1+\varepsilon)\xi)}>
x^{1/(4\xi)}$ and $\tau(m^3)\le L^2$.
The complete earlier counting lemma, with $k=3$, $C=2$, and $c<1$,
therefore bounds the remaining integers by $O_A(xL^{-A})$ for every
fixed $A$. This proves the claimed bound for $B_2$.

Partial summation gives $\sum_{b\in B_2}1/b<\infty$, and hence its
multiples have vanishing upper-density tails. The same holds for $B_1$.
Scaling partitions and taking a least divisor with a partition identify
all strong Mondrian side lengths with the set of multiples of $B$.
Since $1\notin B$, the acquired Behrend-tail argument completes the
proof, including existence of the density.
\end{proof}

\section{Exact source boundary}
This proof recovers the uniform conclusion of Theorem 1.3, not the
literal intermediate estimates of Section 6. The source's displayed
identity $|\tau(\ell^3,\theta)|^2=16\cos^6(\theta\log\ell/2)$
is false: at $\theta\log\ell=\pi/2$ its two sides are zero and two.
Equation~\eqref{polynomial} is the exact replacement.

Also, knowing that a divisor has a prime above the cutoff does not imply
that such a prime occurs with exponent at least two. We make no claim
that the source's particular set $S$ satisfies our stronger assertion.
Lemma~\ref{smallpart} changes the cutoff and proves an additional
negligible-exception condition; Lemma~\ref{assignment} then establishes
the assertion actually needed. In the equal-prime case we use two
prime-counting logarithmic savings, not an unjustified third saving.
In the distinct-prime case the nonsingular exponent matrix is proved
from the canonical assignment, not assumed. The corrected geometric
trace likewise keeps multiplicities on an arbitrary chosen axis.

No Mondrian partition is exhibited or ruled out individually by a
finite search. Conjecture 1.1, the proposed higher-dimensional extension,
and the weaker condition that only the boxes themselves be distinct
remain unanswered. This supplement does not establish the additional
general multiplicative-weight or improved-error assertions in the
source's Section 4 remarks. Those boundaries must remain explicit in
any subsequent whole-source assembly or comparison.

\begin{thebibliography}{9}
\bibitem{LLT} N.~Lebowitz-Lockard and G.~Tenenbaum,
\emph{On solutions to the Mondrian Puzzle}, retained preprint dated
28 August 2026, 13 pages, especially Section 6, pp.~10--12.
\bibitem{prior} Human Compression project,
\emph{Divisor concentration and primitive bases} and
\emph{Primitive rectangular partitions and planar density},
accepted bounded F008 planar packet and its distinct independent review.
\end{thebibliography}
\end{document}
